\documentclass[10pt]{article}

\usepackage[margin=1in]{geometry}
\usepackage[T1]{fontenc}
\usepackage{lmodern}
\usepackage{amsmath,amssymb,amsthm}
\usepackage{booktabs}
\usepackage{enumitem}
\usepackage{graphicx}
\usepackage[hidelinks]{hyperref}
\usepackage[numbers,sort&compress]{natbib}
\usepackage{xcolor}

\newtheorem{definition}{Definition}
\newtheorem{proposition}{Proposition}
\newtheorem{observation}{Observation}
\newtheorem{remark}{Remark}

\newcommand{\E}{\mathbb{E}}
\newcommand{\C}{\mathcal{C}}
\newcommand{\N}{\mathcal{N}}
\newcommand{\Pd}{\mathcal{P}_\delta}
\newcommand{\novel}{\star}
\newcommand{\method}{\textsc{Cira}}

\setlist{itemsep=1pt,topsep=3pt}
\emergencystretch=1.5em

\title{Adapting Before Contact:\\
Compositional Contextual Inference for Safe Control with Frozen World Models}
\author{Research Proposal --- v3}
\date{September 28, 2026}

\begin{document}
\maketitle

\begin{abstract}
Robots that plan with learned world models make mistakes when the physics changes, for example when
the table surface becomes more slippery or the object becomes heavier. Current methods fix this by
updating the model from its prediction errors. They can only do so after the robot has touched the
object, and they discard what they learned when the next episode starts. For tasks that are decided at
the moment of first contact, such as striking a puck toward a target or lifting an unknown load, these
methods therefore give no benefit at all.

We propose a different approach, based on a model of how humans learn motor skills~\cite{heald2021coin}.
The robot keeps a growing set of small correction models, one for each physical condition it has
encountered (``memories''). It continuously estimates which condition it is currently in. It also
learns, without labels, which visual features predict each condition, for example that a grey felt mat
means high friction. When a known condition returns, the robot can select the right memory from what
it sees, before it touches anything.

The method, \method{} (Context-Indexed Residual Adapters), adds these memories to any frozen world
model. Memories for different factors, such as surface and load, can be combined, so a new combination
of known factors can be predicted without new data. The same estimate of the current condition also
drives a safety filter whose safety margins are tuned separately for each memory. The proposal contains two theoretical results and one measurement, each tested on simulated and
real robot tasks:
\begin{itemize}
  \item the error at first contact grows with how unpredictable the sequence of conditions is, while
  for existing methods it grows with the number of episodes in which the condition differs from
  training;
  \item a long-run bound on safety violations that holds whether or not the visual cues are reliable;
  \item a method that splits the robot's recovery after a change into ``recalling an old memory'' and
  ``learning something new'', which we measure over a long real-robot deployment.
\end{itemize}
\end{abstract}

%=====================================================================
\section{Motivation and contributions}
%=====================================================================

\paragraph{World models and their adaptation.}
A latent world model encodes camera images into a compact vector (the \emph{latent state}) and
predicts how that vector changes when the robot acts. Planners use these predictions to choose
actions. Such models now plan manipulation tasks without task-specific
training~\cite{zhou2025dinowm,assran2025vjepa2,sobal2025pldm,hansen2024tdmpc2}. When the physics at
deployment differs from the training data, the predictions become wrong. Two recent methods address
this. AdaJEPA updates part of the model with a few gradient steps on its recent prediction errors while
the robot runs~\cite{wang2026adajepa}. Sandwich-Residuals keeps the model fixed and trains small
correction modules around it~\cite{soni2026sandwich}. Both reset to the original model at the start of
every episode, and both list keeping what was learned across episodes as future work. Keeping the
gradient updates running indefinitely is not a good alternative. On long streams of data such updates
gradually drift and degrade, and simply resetting from time to time works better~\cite{press2023rdumb,hoang2024petta}.

\paragraph{Real deployments repeat themselves.}
A robot in a lab, warehouse or kitchen rarely meets a completely new condition. It meets the same few
surfaces, containers and tool states repeatedly, and these usually look different from each other.
Humans make use of this. Research on human motor learning shows that people keep separate motor
memories for different conditions. They recall a memory when a condition returns instead of
relearning it (\emph{savings}). Learning a new condition does not erase the old one. Visual cues can
select the right memory before the movement starts~\cite{heald2021coin,heald2023annurev,haruno2001mosaic}.

\paragraph{Running example.}
A robot arm strikes a puck toward a target on a table. The table is covered with one of three mats:
felt (high friction), rubber (medium) or acrylic (low). The puck is sometimes empty and sometimes
weighted. The mats and weights are swapped by people during the day. If the robot assumes the wrong
friction, the puck stops short or slides off the table edge. The strike is a single contact, so the
robot has no chance to correct its action after seeing the error.

\paragraph{Main idea.}
We treat adaptation as the question ``which of the conditions I know is active now, or is this a new
one?'' rather than ``how should I change my model parameters?''. The robot answers this question with
a probability distribution over its memories (the \emph{context posterior}). It updates only the
memory that is likely to be active. This matters most at the first contact. At that point
error-driven methods have received no error signal. The outcome of single-contact tasks is decided
there, and so is much of the risk of unsafe behaviour.

\paragraph{Contributions.}
\begin{enumerate}[label=\textbf{C\arabic*.},leftmargin=2.2em]
  \item \textbf{Single-contact tasks as a test.} We define a class of tasks in which methods that only
  learn from prediction errors cannot outperform the unadapted model (Obs.~\ref{obs:reactive}).
  Examples are striking an object toward a target, lifting an unknown load and pushing near a table
  edge. On these tasks, only information available before contact can help.
  \item \textbf{The method, \method{}.} A set of small probabilistic correction models (memories) on
  top of a frozen world model, with closed-form updates. One probability distribution over memories
  decides when a new memory is created, which memory is used and updated, and when the robot has
  returned to normal conditions. A noise model that depends on the state prevents the short, noisy
  moments at contact from being mistaken for a change of condition. Memories are organised by factor
  (for example surface and load), each with its own visual cue. As a result, an unseen combination of
  known factors can be predicted before contact.
  \item \textbf{Theory.} Proposition~\ref{prop:separation} bounds the total error at first contact.
  For our method the bound depends on how predictable the sequence of conditions is given the visual
  cues. For reset-based methods it grows with every episode that differs from training.
  Proposition~\ref{prop:latency} predicts how many steps the robot needs to identify a known
  condition. We state these predictions before running the experiments.
  \item \textbf{A safety filter that uses the memories.} The filter checks each planned action against
  every memory that is still reasonably likely. Its safety margin is tuned online, separately for each
  memory. When a known condition returns, the robot therefore also recovers the margin that was
  appropriate for it. Proposition~\ref{prop:safety} bounds the long-run rate of safety violations for
  any sequence of conditions, including when visual cues are misleading.
  \item \textbf{A measurement and a long deployment.} We split the improvement after a change of
  condition into two parts (Def.~\ref{def:decomp}): the part that comes from selecting an existing
  memory, and the part that comes from changing a memory's parameters. Following the human motor
  learning literature, we call these \emph{apparent} and \emph{proper} learning. We report this split
  over a real-robot deployment of at least 1000 episodes, in which people swap surfaces and containers
  on a schedule the robot does not know. We also release a benchmark of long sequences of changing
  conditions, built around four classic motor-learning experiments.
\end{enumerate}

\paragraph{Scope.}
We study changes in physics (friction, mass, load, motor strength) while the image encoder stays fixed.
Appearance matters only as a clue about which condition is active. We do not try to correct for
changes in appearance that confuse the encoder itself, which is what AdaJEPA and Sandwich-Residuals
address. We also assume that a safety constraint expressed in the latent space matches the real
constraint. Whether it does is a property of the encoder that we do not study here~\cite{lutkus2025latent}.

%=====================================================================
\section{Related work}
%=====================================================================

\paragraph{Contextual inference in human motor learning.}
The MOSAIC model proposed that the brain holds several paired prediction and control models and
chooses between them using both prediction accuracy and signals available before a
movement~\cite{wolpert1998multiple,haruno2001mosaic}. The COIN model~\cite{heald2021coin} turned this
into a full probabilistic model. It has the following parts:
\begin{itemize}
  \item a set of contexts that can grow without limit, which tend to persist over time;
  \item one simple memory per context;
  \item visual or other cues whose probability depends on the context.
\end{itemize}
COIN explains several well-known effects by changes in which memory is used, not by changes in
learning speed. These are savings, interference between opposing tasks, and the spontaneous return of
old behaviour. It also predicted two effects that the standard two-rate model~\cite{smith2006interacting}
cannot produce, and experiments confirmed both. One is the return of an old memory triggered by a few
reminder trials (\emph{evoked recovery}). The other is that the amount learned in one trial depends on
the cue. Related models of animal learning make the same distinction between updating a memory and
creating a new one~\cite{gershman2010context,gershman2014statistical,oh2019minimizing}. People also
learn separate parts of a task independently and recombine them~\cite{franklin2018compositional}.

One result is a warning for our design. Many purely visual cues do not help people separate motor
memories~\cite{howard2013contextual}, while differences in the planned movement do~\cite{sheahan2016planning}.
We therefore do not assume that a cue is useful. The robot learns how informative each cue is, and we
vary cue reliability in the experiments.

\paragraph{Adapting world models while they run.}
AdaJEPA~\cite{wang2026adajepa} updates encoder and predictor layers with about one gradient step per
planning cycle on the PushT, PushObj and PointMaze tasks, including changes in mass and damping. Its
authors note that adaptation is limited by what the pretrained representation can express.
Sandwich-Residuals~\cite{soni2026sandwich} reaches similar results while training only 1--3\% as many
parameters. More general test-time adaptation methods~\cite{sun2020ttt,wang2021tent,wang2022cotta}
degrade on long and repeating data streams~\cite{press2023rdumb,hoang2024petta}. Contextual world
models estimate a continuous context vector from recent
experience~\cite{roder2025dali,benad2026dmash,lee2020context}. Changepoint-aware world models
discard outdated training data after a detected change~\cite{yang2026cawm}. None of these methods
stores separate conditions and recalls them later, and none acts before contact.

\paragraph{Probabilistic tools for switching dynamics.}
Our method combines existing tools:
\begin{itemize}
  \item Bayesian last layers, which give fast closed-form updates of the final layer of a
  network~\cite{harrison2018alpaca,harrison2024vbll};
  \item methods that detect when the data-generating process
  changes~\cite{harrison2020moca,adams2007bocpd};
  \item mixtures of dynamics models whose number grows with the data, using a Chinese restaurant process
  (CRP) prior~\cite{nagabandi2019mole,xu2020taskagnostic,jerfel2019reconciling,lee2020cndpm};
  \item switching linear dynamical systems~\cite{fox2011bnpslds,fox2011sticky,linderman2017rslds};
  \item the interacting multiple-model (IMM) filter from target
  tracking~\cite{blom1988imm,mazor1998imm,li1996vsmm};
  \item lifelong learning with factored or modular
  models~\cite{feng2022factored,mendez2022modular,killian2017robust}.
\end{itemize}
The closest prior method is MOLe~\cite{nagabandi2019mole}, which keeps a CRP mixture of dynamics models
for model-based reinforcement learning. It selects a model only from prediction errors, so it cannot
act before contact. In online learning theory, ``mixing past posteriors''~\cite{bousquet2002tracking}
handles a small set of recurring experts. Our first-contact result builds on the same intuition.

\paragraph{Using vision to predict dynamics in advance.}
VcVBLL~\cite{brunzema2026vcvbll} adjusts a vehicle dynamics model using a hand-designed score for
standing water in the camera image. Terrain-aware models use pretrained image features in a similar
way~\cite{gibson2024visualterrain,cai2023probtrav,agarwal2022egocentric}. Other methods train an image
model online to predict physical properties measured by the robot's own
sensors~\cite{loquercio2022crossmodal,margolis2023activesensing,frey2023wvn}. Vision-language models
can also provide prior estimates of physical properties~\cite{gao2024physically,wang2026phys2real}. All
of these map an image to a continuous physical quantity through a channel chosen in advance. \method{}
instead maps the image to one of the robot's own discovered conditions, and lets the observed dynamics
overrule the image when they disagree.

\paragraph{Safety in latent space.}
Several methods compute safety filters directly in the latent space of a world model, using
reachability analysis or control barrier functions
(CBFs)~\cite{nakamura2025latent,nakamura2026latentcbf,anand2025safety,lutkus2025latent,wu2026points}.
A CBF is a function $h$ that is positive in safe states. A controller stays safe if it keeps $h$ from
decreasing too quickly. Seo et al.~\cite{seo2025uncertainty} treat states outside the training data as
unsafe and steer away from them. AnySafe changes the safety constraint at runtime but not the
dynamics~\cite{agrawal2025anysafe}. Conformal prediction gives error bounds for learned predictors,
and adaptive conformal inference (ACI) keeps these bounds valid when the data change over
time~\cite{gibbs2021adaptive,lindemann2023safe,dixit2023adaptive,huriot2026safe,nath2026pixels,srinivasan2026safety}.
Adaptive and robust CBFs narrow down a single set of unknown parameters
online~\cite{taylor2020adaptive,lopez2021robust,taylor2020learning,castaneda2025recursively}.
Multi-stage model predictive control is robust against a fixed set of candidate
models~\cite{lucia2013multistage}. What is missing is a filter whose set of candidate models is the
set of remembered conditions that are currently likely. Such a filter can be cautious in new
conditions, less cautious in recognised ones, and can store its calibration with each memory.

\begin{table}[t]
\centering
\footnotesize
\setlength{\tabcolsep}{2.5pt}
\resizebox{\linewidth}{!}{%
\begin{tabular}{lccccccc}
\toprule
 & Latent visual & Keeps / & Acts before & Cue learned & Combines & Safety uses & Calibrated \\
 & world model & recalls memories & contact & during use & factors & memories & under change \\
\midrule
AdaJEPA~\cite{wang2026adajepa}, Sandwich-Res.~\cite{soni2026sandwich} & \checkmark & -- & -- & -- & -- & -- & -- \\
MOLe~\cite{nagabandi2019mole}, IMGP~\cite{xu2020taskagnostic} & -- & \checkmark & -- & -- & -- & -- & -- \\
MOCA~\cite{harrison2020moca} & -- & keeps only & -- & -- & -- & -- & partly \\
VcVBLL~\cite{brunzema2026vcvbll}, Phys2Real~\cite{wang2026phys2real} & -- & -- & \checkmark & -- (fixed in advance) & -- & -- & \checkmark \\
Uncertainty-aware latent filter~\cite{seo2025uncertainty} & \checkmark & -- & -- & -- & -- & avoids unknown & \checkmark \\
Latent world model + ACI filters~\cite{huriot2026safe,nath2026pixels} & \checkmark & -- & -- & -- & -- & one model & \checkmark \\
COIN~\cite{heald2021coin} (model of humans) & -- & \checkmark & \checkmark & \checkmark & -- & -- & -- \\
\midrule
\method{} (this proposal) & \checkmark & \checkmark & \checkmark & \checkmark & \checkmark & \checkmark & \checkmark \\
\bottomrule
\end{tabular}}
\caption{Comparison with the closest work, based on literature searches in September 2026. Each
column is covered by some prior work. Our contribution is the combination, the single-contact task
class and the first-contact result.}
\label{tab:positioning}
\end{table}

%=====================================================================
\section{Problem setting}
%=====================================================================

A frozen encoder maps an image $o_t$ to a latent state $z_t=E(o_t)\in\mathbb{R}^d$. A frozen predictor
gives the next latent state $\hat z_{t+1}=f(z_t,a_t)$ for action $a_t$. The robot runs one long
sequence of episodes. It knows when an episode starts and ends, but it is never told which condition
is active, and its memories are never reset.

In episode $e$ an unknown \emph{context} $c^\ast_e$ is active. A context is a physical condition, such
as ``acrylic mat, weighted puck''. Under context $c$, the true next state differs from the frozen
prediction by an average amount $\mu_c(z,a)$, which we call the \emph{residual}. The training
condition has zero residual, $\mu_0\equiv 0$. Before the first contact, at time $t_c$, the robot sees
the scene and extracts a visual \emph{cue} $y_e$. Before contact, the robot's own motion reveals
nothing about the context.

\begin{definition}[Single-contact task]
A task is \emph{single-contact} if its success, or a violation of its safety constraint, is decided by
the prediction the model makes at the first contact $t_c$ and the action chosen from it. Examples are
striking a puck toward a target, tossing an object~\cite{zeng2020tossingbot}, other fast
non-grasping motions~\cite{lynch1999dynamic}, lifting an unknown load without jerking or slipping, and
pushing an object near a table edge on an unknown surface.
\end{definition}

\begin{observation}[Error-driven adaptation cannot help at first contact]\label{obs:reactive}
Suppose a method changes its model only in response to prediction errors. Since there are no errors
before contact, its prediction at $t_c$ is the same as at the start of the episode. Methods that reset
every episode (AdaJEPA, Sandwich-Residuals, gradient updates with reset) therefore make the frozen
model's prediction at first contact. Their error there is $\|\mu_{c^\ast_e}\|$. On single-contact tasks
they cannot do better than the unadapted model.
\end{observation}

The observation follows directly from the definitions. Its purpose is to make clear that on these
tasks only two sources of information remain: what the robot has seen in past episodes, and what it
sees in the scene before contact.

%=====================================================================
\section{Method: Context-Indexed Residual Adapters}
%=====================================================================

\paragraph{Overview.}
\method{} has four parts:
\begin{itemize}
  \item a small set of directions in the latent space where corrections are made;
  \item one correction model (memory) per known context;
  \item a probability distribution over which memory is active, updated at every step;
  \item a learned link between visual cues and memories.
\end{itemize}
The world model itself is never changed.

\paragraph{Where corrections are made.}
We collect the prediction errors $r_t=z_{t+1}-f(z_t,a_t)$ of the frozen model on training-condition
data. We keep the $k$ directions along which these errors vary most (principal components, with
$k$ between 4 and 16), stored as a matrix $U\in\mathbb{R}^{d\times k}$. All corrections are made in
these $k$ directions, using the projected error $e_t=U^\top r_t$. The reasoning is that a few physical
parameters, such as friction and mass, change the dynamics along only a few
directions~\cite{doshivelez2016hipmdp,killian2017robust}.

\paragraph{Memories.}
Each context $c$ has a memory: a linear model that predicts the projected error from a small feature
vector $\phi_t=\phi(U^\top z_t,a_t)$, for example $\phi_t=[a_t,\,U^\top z_t,\,1]$. The weights $W_c$ are
uncertain, and we keep a Gaussian distribution over them:
\begin{equation}
  e_t\mid c\;\sim\;\N\big(W_c\phi_t,\;\Sigma(z_t,a_t)\big),\qquad
  W_c\sim\mathcal{MN}(M_c,\Sigma_{\text{row}},V_c).
  \label{eq:adapter}
\end{equation}
Here $\mathcal{MN}$ is a matrix normal distribution. It is the matrix version of a Gaussian, so the
memory can be updated in closed form after every step~\cite{harrison2018alpaca}. The model includes
the action, because a change in friction or motor strength produces an error that grows with the
action. A constant offset could not describe this. The correction is also linear in the action, which
the safety filter in Section~\ref{sec:safety} uses. The weights $W_c$ are allowed to drift slowly over
time, so each memory can keep learning while it is in use. The memory for the training condition is
fixed at $W_0=0$. One extra ``new context'' slot, written $\novel$, starts from a broad prior and is
used when no stored memory fits.

\paragraph{Noise that depends on the state.}
The noise covariance $\Sigma(z,a)$ comes from a small network trained on training-condition
errors~\cite{kersting2007hetgp}. Around contact the frozen model is already less accurate, for
example at the moment of impact or during stick--slip motion, and $\Sigma$ is larger there. Large
errors at those moments are then treated as expected noise rather than as evidence of a new condition.
Pushing is known to have this kind of state-dependent noise and variable
friction~\cite{bauza2017probpush,ma2018friction}. We found no prior work that links contact noise to
the false creation of new contexts. We expect it to be the most common failure of a simple
implementation, and our simulated tasks include it on purpose.

\paragraph{Combining factors.}
A context can consist of several independent factors, such as the surface and the load:
$c=(c^1,\dots,c^F)$. Each factor has its own set of memories $W^f_{c^f}$, and the corrections add up:
$e_t\mid c\sim\N\big(\sum_f W^f_{c^f}\phi_t,\Sigma\big)$. Two things let the robot separate the factors
without labels:
\begin{itemize}
  \item \emph{Each factor has its own image region.} The surface cue comes from the table area around
  the object. The load cue comes from a crop of the object or container.
  \item \emph{Factors change at different times.} People swap mats and containers independently, so
  the robot sees one factor change while the other stays the same.
\end{itemize}
Adding the corrections is an approximation. Friction force grows with the normal force, so mass and
friction interact. We therefore also keep a memory for each combination the robot has actually seen.
For a combination it has not seen, it starts from the sum of the factor memories and refines this
memory as data arrive. Our claim about new combinations concerns the prediction at their first
contact.

\paragraph{Estimating the active context.}
At every step the robot updates a probability distribution $\pi_t$ over the stored memories, the
training condition and the ``new context'' slot. Three sources of evidence are combined:
\begin{itemize}
  \item \emph{How often each context has occurred.} Before an episode, a context that has been seen
  often is considered more likely. This is the Chinese restaurant process
  prior~\cite{neal2000dpm}. It gives each known context a probability proportional to how often it
  has occurred, and gives a new context a small fixed probability.
  \item \emph{The visual cue}, through the cue model described below.
  \item \emph{How well each memory predicts the latest step}, through Eq.~\eqref{eq:adapter}.
\end{itemize}
Within an episode, contexts are assumed to persist, with a small chance of
switching~\cite{fox2011sticky}. The update
follows the interacting multiple-model filter~\cite{blom1988imm}:
\[
  \pi_t(c)\;\propto\;p(e_t\mid c,\phi_t)\;\prod_f p(y^f\mid c^f)\;\sum_{c'}p(c\mid c')\,\pi_{t-1}(c').
\]
This one distribution answers four questions, so no separate change detector is needed:
\begin{itemize}
  \item if the ``new context'' slot stays most likely for several steps, a new memory is created;
  \item if a stored memory is most likely, it is recalled;
  \item if the training condition is most likely, the robot has returned to normal;
  \item if the distribution shifts, the condition has changed.
\end{itemize}
Each memory is updated in proportion to its probability, as in COIN~\cite{heald2021coin}, rather than
only the most likely one. Two memories that become nearly identical are merged.

\paragraph{Learning what cues mean.}
For each factor value, the robot keeps a Gaussian model $p(y^f\mid c^f)$ of what the cue looks like
under that context. After each episode, this model is updated using context probabilities computed
\emph{from the dynamics only}, without the cue. If the cue model were instead trained on decisions that the cue itself had
influenced, a wrong cue could reinforce itself. With this rule, a misleading cue is gradually
unlearned.

\paragraph{Optional: priors from vision-language models.}
A vision-language model that estimates physical properties~\cite{gao2024physically,wang2026phys2real}
could set the starting guess for a new context from the image. This would help even the first time a
condition is seen. We treat this as an extension and keep it separate from the main claims.

%=====================================================================
\section{Planning and safety in latent space}\label{sec:safety}
%=====================================================================

\paragraph{Planning.}
We use a sampling-based planner, the cross-entropy method (CEM) or MPPI~\cite{rubinstein1999cem,chua2018pets}.
Each sampled action sequence is simulated with the frozen model plus a correction,
$\hat z_{t+1}=f(z_t,a_t)+UW\phi_t$. The weights $W$ come from a context drawn from $\pi_t$. The planner
can minimise the average cost over contexts, or a risk-averse version (CVaR) that focuses on the worst
outcomes~\cite{yin2023risk,lew2024risk}. Updating $\pi_t$ and the memories costs much less than one
call of the world model.

\paragraph{Which contexts the safety filter considers.}
The planner uses the full distribution $\pi_t$. The safety filter uses only the contexts with
probability at least $\delta$, the \emph{plausible set} $\Pd=\{c:\pi_t(c)\ge\delta\}$. This set includes
the ``new context'' slot whenever that slot's probability exceeds $\delta$. When the robot recognises
a known condition, $\Pd$ contains a single memory. When it is unsure, or the condition looks new, $\Pd$
contains several, and the filter becomes more cautious. As a baseline we also consider a filter that
guards against \emph{all} stored memories at once, as in multi-stage model predictive
control~\cite{lucia2013multistage}. We expect that baseline to be much more cautious.

\paragraph{The safety condition.}
We assume a latent safety function $h(z)$ is given, for example a latent CBF or reachability value
trained with existing methods~\cite{nakamura2025latent,nakamura2026latentcbf}. States with $h(z)\ge 0$
are safe. The filter accepts an action $a$ only if the discrete-time CBF
condition~\cite{agrawal2017discrete} holds for every plausible memory:
\begin{equation}
  h\big(f(z_t,a)+UW_c\phi(z_t,a)\big)\;-\;L_h\,q^{c}_t\;\ge\;(1-\alpha)\,h(z_t)
  \qquad\text{for all } c\in\Pd .
  \label{eq:cbf}
\end{equation}
In words, $h$ must not drop by more than a fraction $\alpha$ in one step under any plausible memory,
with an extra safety margin $L_h q^c_t$. Here $L_h$ bounds how fast $h$ can change, and $q^c_t$ is an
estimate of how wrong memory $c$ can be, explained below. Because the correction is linear in the
action, linearising $f$ and $h$ around a planned action turns the condition into one linear constraint
per plausible memory. The filter is then a small quadratic program. Over the planning horizon we
apply the same condition by discarding sampled action sequences that violate it. If no action
satisfies it, the robot stops or retracts.

\paragraph{Safety margins stored with each memory.}
After each step we measure how far the actual latent state was from the nearest plausible prediction:
$s_t=\min_{c\in\Pd}\|r_t-UW_c\phi_t\|$. This distance includes ordinary noise, errors outside the $k$
correction directions and mistakes in identifying the context. The margin $q^c_t$ is set by adaptive
conformal inference (ACI)~\cite{gibbs2021adaptive}. ACI raises the margin after each step where the
distance exceeded it and lowers it slightly otherwise, so that in the long run the margin is exceeded
on a chosen fraction $\delta_{\mathrm s}$ of steps. We run a separate ACI for each memory and update it
only on steps attributed to that memory. When a known condition returns, its margin is therefore
already tuned and does not have to be learned again.

\paragraph{What happens when the cue is wrong.}
The cue affects $\pi_t$ and therefore $\Pd$. A reliable cue shrinks the plausible set before contact,
so the robot can act less cautiously from the start. A misleading cue can leave the true context out
of $\Pd$. In that case the measured distances grow and ACI raises the margin. The long-run bound in
Proposition~\ref{prop:safety} still holds. A wrong cue therefore costs extra caution and a limited
number of extra violations. We measure this cost against a filter that ignores cues.

\paragraph{Optional: safe probing.}
When the robot is unsure which context is active, it can prefer actions that help tell contexts
apart. Among the candidate actions that satisfy Eq.~\eqref{eq:cbf} and whose task cost is within a
factor $(1+\epsilon)$ of the best, CEM favours those with the highest expected information about the
context, $I(c;e_{t+1:t+H}\mid\mathbf a)$. This idea comes from active fault
diagnosis~\cite{blackmore2006discrimination,blackmore2008active,simandl2009afd} and dual
control~\cite{mesbah2018dual,heirung2019mpcdiscrimination}, combined with safe
exploration~\cite{koller2018learning,berkenkamp2017safe}. In a single-contact task this means the robot
gives the puck a small nudge first, but only when the cue and its past experience leave the context
unclear. How often it chooses to probe is itself a result we report.

%=====================================================================
\section{Theory}
%=====================================================================

\begin{proposition}[Error at first contact]\label{prop:separation}
Consider $N$ episodes whose contexts come from a finite set of $K$ contexts, visual cues that take
values in a finite set $\mathcal{Y}$, and memories whose mean predictions are accurate to within
$\varepsilon$. Let $\Delta_{\max}=\max_{c,c'}\|\mu_c-\mu_{c'}\|$ be the largest difference between two
contexts, and $\Delta_{\min}=\min_{c\ne0}\|\mu_c\|$ the smallest difference between any context and the
training condition.
\begin{enumerate}[label=(\roman*)]
  \item Any method that resets every episode has total first-contact error
  $\sum_e\|\mu_{c^\ast_e}\|\ge N_{\mathrm{shift}}\Delta_{\min}$. Here $N_{\mathrm{shift}}$ is the number
  of episodes whose context differs from training.
  \item In episode $e$, \method{} has first-contact error at most
  $\varepsilon+(1-\pi_0^e(c^\ast_e))\Delta_{\max}$, where $\pi_0^e(c^\ast_e)$ is the probability it
  assigns to the true context before contact. Summed over episodes,
  $\sum_e(1-\pi_0^e(c^\ast_e))\le\sum_e-\log\pi_0^e(c^\ast_e)$. The right-hand side is the total
  log-loss of predicting each episode's context from its cue. For any sequence of contexts it is at
  most $N\hat H(C\mid Y)+O(K|\mathcal{Y}|\log N)$~\cite{krichevsky1981performance}, where
  $\hat H(C\mid Y)$ is the empirical conditional entropy of the context given the cue.
\end{enumerate}
\method{} therefore has lower total first-contact error than any reset-based method whenever
$\varepsilon+\Delta_{\max}\hat H(C\mid Y)<\Delta_{\min}N_{\mathrm{shift}}/N$ (entropy in nats). This
ignores lower-order terms and the $O(K)$ episodes in which memories are first learned. It holds for
any order of contexts, including adversarial ones.
\end{proposition}

\noindent\emph{Proof sketch.} Part (i) follows from Observation~\ref{obs:reactive}. For (ii), the
error of the probability-weighted prediction is
$\|\sum_{c\ne c^\ast}\pi(c)(\mu_c-\mu_{c^\ast})\|\le(1-\pi(c^\ast))\Delta_{\max}$, and $1-x\le-\log x$
for $x\in(0,1]$. The log-loss bound is the standard bound for sequential prediction with a Dirichlet
mixture, applied separately to each cue value. Creating new contexts adds $O(K\log N)$.

\paragraph{Interpretation.}
The conditional entropy $\hat H(C\mid Y)$ measures how hard it is to guess the context of an episode
from its cue and from past frequencies. If conditions repeat and the cue identifies them, it is close
to zero and \method{} makes almost no first-contact error. If every episode brings a new condition,
memory cannot help and the bound says so. The benefit of the cue is the reduction
$\hat H(C)-\hat H(C\mid Y)$, which is the information the cue carries about the context. The same
quantity therefore explains three effects: savings, correct action before contact, and the lack of
benefit from an uninformative cue.

\begin{proposition}[Steps needed to recognise a known context]\label{prop:latency}
Consider an episode whose true context $c^\ast$ is stored in memory, and let $c'$ be the most likely
alternative. Suppose that after contact each step adds, on average, $D>0$ to the log-odds in favour of
$c^\ast$, where $D$ is the KL divergence between the two memories' predictions, and that the step
contributions are independent. Before contact they are zero. The starting log-odds combine past
frequency and cue:
$\lambda_0=\log\frac{n_{c^\ast}}{n_{c'}}+\log\frac{p(y\mid c^\ast)}{p(y\mid c')}$, where $n_c$ counts
past occurrences of $c$. If $\lambda_0$ already exceeds the decision threshold $\lambda^\ast$, the
number of steps needed is zero. Otherwise, by Wald's identity, the expected number of steps is at least
$(\lambda^\ast-\lambda_0)/D$, with equality apart from the final overshoot.
\end{proposition}

This gives predictions we can write down before running experiments:
\begin{itemize}
  \item \textbf{Acting before contact.} A learned cue that is correct with probability $\rho$ lets
  the robot correct about a fraction $\rho$ of the error before contact.
  \item \textbf{Savings.} Each further occurrence of a context shortens recognition by
  $\log(n_{k+1}/n_k)/D$ steps, so the gains shrink over time. For gradient updates with reset, the time
  stays the same.
  \item \textbf{Misleading cues.} A misleading cue of strength $\ell_y$ (its log-likelihood ratio)
  costs about $2\ell_y/D$ extra steps compared with a correct cue.
  \item \textbf{Noise.} The state-dependent noise model lowers $D$ around contact. This makes
  recognition slower but prevents false new contexts. We measure both effects. For the case without
  cues, classical results on detection delay~\cite{lorden1971procedures,lai1998information} serve as
  a reference.
\end{itemize}

\begin{proposition}[Long-run safety for any cue quality]\label{prop:safety}
Suppose Eq.~\eqref{eq:cbf} is enforced exactly whenever some action satisfies it (for example by
rejecting sampled actions), and the fallback action is safe otherwise. Suppose ACI runs separately
for each memory with target rate $\delta_{\mathrm{s}}$ and step size $\gamma$. Then, for any sequence
of contexts and cues, the one-step condition $h(z_{t+1})\ge(1-\alpha)h(z_t)$ fails on at most a
fraction $\delta_{\mathrm{s}}+O\big(K/(\gamma T)\big)$ of the first $T$ steps.
\end{proposition}

\noindent\emph{Proof sketch.} If the measured distance $s_t$ is within the margin $q^c_t$, the true
next state lies within $q^c_t$ of the prediction of some plausible memory. Eq.~\eqref{eq:cbf} and the
bound $L_h$ on how fast $h$ changes then give the one-step condition. So the condition can fail only on
steps where the margin was exceeded. For each memory, ACI guarantees that the long-run fraction of
such steps is at most $\delta_{\mathrm{s}}+O(1/(\gamma T_c))$, without assumptions on the data
distribution~\cite{gibbs2021adaptive}. Summing over the $K$ memories gives the result.

The bound applies to the safety condition in latent space. Whether this implies physical safety
depends on whether $h$ correctly represents the real constraint, which we treat as a property of the
encoder~\cite{lutkus2025latent}. In single-contact tasks it is natural to count episodes instead of
steps, and the bound becomes a long-run rate of unsafe episodes.

%=====================================================================
\section{Measuring recall versus learning}
%=====================================================================

For every adaptive method in this proposal, we separate its adaptive state into two parts:
\begin{itemize}
  \item \emph{which memories it uses}, $\pi$: for \method{} this is the context distribution;
  \item \emph{what the memories contain}, $\Theta$: the weights.
\end{itemize}
Methods that reset every episode have only $\Theta$. Let $L(\pi,\Theta)$ be the prediction error on
data from the current condition, computed by replaying recorded transitions.

\begin{definition}[Splitting recovery into recall and learning]\label{def:decomp}
After a change at time $t_0$, the improvement by time $t$ is
$R_t=L(\pi_{t_0},\Theta_{t_0})-L(\pi_t,\Theta_t)$. The part due to changing which memories are used is
\[
  A_t=\tfrac12\big[L(\pi_{t_0},\Theta_{t_0})-L(\pi_t,\Theta_{t_0})\big]
     +\tfrac12\big[L(\pi_{t_0},\Theta_t)-L(\pi_t,\Theta_t)\big],
\]
the average effect of switching $\pi$ with $\Theta$ held at its old and at its new value (a Shapley
value with two players). We call $A_t$ \emph{apparent learning}. The remainder $P_t=R_t-A_t$ is
\emph{proper learning}.
\end{definition}

The two parts always add up to the total, and the result does not depend on which part is computed
first. For \method{} every term can be computed in closed form from stored distributions. For methods
that reset every episode, $A_t$ is always zero. We apply the same split to the safety margin. This
shows how much of the reduction in caution after a change comes from recalling a memory and how much
from recalibration.

%=====================================================================
\section{Hypotheses and stopping criteria}
%=====================================================================

\begin{enumerate}[label=\textbf{H\arabic*},leftmargin=2.2em]
  \item \textbf{First contact.} On single-contact tasks with recurring conditions and an informative
  cue, \method{} succeeds more often than the frozen model and all error-driven baselines. Its
  first-contact error follows Proposition~\ref{prop:separation}.
  \item \textbf{Savings.} Recognition time falls with each further occurrence, as predicted by
  Proposition~\ref{prop:latency}. For gradient updates with reset it does not change.
  \item \textbf{No interference.} Learning condition B does not reduce performance on condition A.
  Continuously running gradient updates and continual test-time adaptation do reduce it.
  \item \textbf{Misleading cues.} The observed dynamics overrule a misleading cue within the predicted
  number of steps. Task cost stays below that of the frozen model, and the safety bound holds.
  \item \textbf{New combinations.} For a combination of factors never seen together (for example the
  acrylic mat with the weighted puck), the combined factor memories give lower first-contact error than
  both the version without cues and the version without factors.
  \item \textbf{Safety.} The long-run rate of violations of the latent safety condition stays within
  $\delta_{\mathrm s}$ plus the predicted slack, also with misleading cues. At the same violation rate,
  our filter intervenes less than a filter that guards against all memories and less than a filter
  with one shared margin. Its caution decreases as conditions recur.
  \item \textbf{Noise.} With noisy contacts, no single learning rate for gradient updates is best both
  in normal conditions and after a change. \method{} matches the best rate in each phase. Its rate of
  falsely created contexts stays below a limit fixed in advance.
  \item \textbf{Recall versus learning.} In the long real deployment, recall accounts for a substantial
  share of recovery. Fixed in advance: the median of $A/R$ is at least 0.3.
\end{enumerate}

\paragraph{When we change the plan.}
\begin{itemize}
  \item If the version without cues is as good as the full method at first contact, we drop the claim
  about acting before contact. The paper then focuses on memory, the benchmark and the recall/learning
  measurement.
  \item If well-tuned gradient updates with reset are as good as the version without cues on recurring
  conditions, the method offers no advantage. The paper then reports the benchmark and this negative
  result.
  \item If combining factors fails on unseen combinations, we report whole-context memories only and
  drop the claim about new combinations.
  \item If our safety filter is not less cautious than the all-memories filter at the same violation
  rate, we present the safety part as a supporting result rather than a main contribution.
\end{itemize}

%=====================================================================
\section{Experiments}
%=====================================================================

\paragraph{Sequences of conditions.}
Each condition lasts for a random number of episodes (geometric distribution), and it occasionally
changes within an episode. Each condition has a matching appearance. The cue shows that appearance
with probability $\rho$ and a different one otherwise. We test $\rho\in\{1/M,0.5,0.65,0.8,0.95\}$,
where $M$ is the number of conditions and $\rho=1/M$ means the cue carries no information. The
sequences contain four patterns taken from motor-learning experiments:
\begin{itemize}
  \item \emph{savings}: conditions A, B, A, B in blocks;
  \item \emph{interference}: a long or a short block of A before B;
  \item \emph{evoked recovery}: two episodes of A inserted into a block of normal conditions;
  \item \emph{misleading cues}: the pairing between cue and condition swapped for one block.
\end{itemize}
Each pattern is linked to one of the predictions above. A pattern that does not separate the methods
in simulation is removed.

\paragraph{Environments.}
\begin{enumerate}[label=(\roman*)]
  \item \textbf{Simulation.} A 2D pushing task and a 2D puck strike with a table edge. Contacts are
  noisy, and surface and load vary. Cues are correct 80\% of the time. These tasks run on a laptop
  and decide whether H1, H2, H4, H5 and H6 hold before we invest in hardware.
  \item \textbf{AdaJEPA tasks}~\cite{wang2026adajepa}. PushT, PushObj and PointMaze with the released
  JEPA and DINO-WM~\cite{zhou2025dinowm} models. We pair the existing colour changes of the agent or
  block with changes in dynamics, so that colour becomes a cue, and we add random noise to the
  dynamics. This allows a direct comparison with AdaJEPA and Sandwich-Residuals.
  \item \textbf{Real robot, single-contact tasks.}
  \begin{itemize}
    \item \emph{Strike:} strike a puck to a target on felt, rubber or acrylic mats, without it
    sliding off the table edge~\cite{yu2016million}.
    \item \emph{Lift and place:} lift loads in visually distinct containers without exceeding a limit
    on jerk or tilt.
  \end{itemize}
  The world model is a frozen V-JEPA~2-AC~\cite{assran2025vjepa2} or a DINO-WM-style model trained on
  training-condition data. The memories add only a few kilobytes on top of a large pretrained model
  that stays unchanged.
  \item \textbf{Real robot, long deployment.} At least 1000 episodes over several days. People swap
  mats and containers according to a schedule the robot does not know, and sometimes cover the mats
  with a neutral sheet to make cues less reliable. We report the recall/learning split, predicted
  versus observed first-contact error, and the violation rate.
\end{enumerate}

\paragraph{Baselines.}
\begin{itemize}
  \item Adaptation:
  \begin{itemize}
    \item the frozen model;
    \item continuously running gradient updates (as in AdaJEPA);
    \item gradient updates with periodic reset~\cite{press2023rdumb};
    \item gradient updates reset exactly at each change (an idealised version);
    \item Sandwich-Residuals~\cite{soni2026sandwich};
    \item a MOLe-style mixture of models~\cite{nagabandi2019mole};
    \item MOCA-style change detection without recall~\cite{harrison2020moca}.
  \end{itemize}
  \item Safety:
  \begin{itemize}
    \item the uncertainty-aware latent filter~\cite{seo2025uncertainty};
    \item a latent filter with one shared ACI margin~\cite{huriot2026safe};
    \item a filter that guards against all memories~\cite{lucia2013multistage}.
  \end{itemize}
  \item Ablations of \method{}:
  \begin{itemize}
    \item without cues;
    \item without factors;
    \item with constant noise;
    \item with a constant-offset memory;
    \item updating only the most likely memory;
    \item with one shared margin instead of one per memory.
  \end{itemize}
  \item Upper bound: \method{} told the true context.
\end{itemize}
All settings are tuned on training-condition data only. Random seeds and sequences are fixed and
published before the main experiments.

\paragraph{Metrics.}
\begin{itemize}
  \item Error at first contact, and success on single-contact tasks.
  \item Number of steps after contact needed to recognise the condition.
  \item Total extra prediction error after each change.
  \item Savings: recognition time for the $k$-th occurrence relative to the first.
  \item Performance on condition A after learning condition B.
  \item Error on unseen combinations of factors.
  \item Violation rate, and how much the safety filter changes the planned action or raises task cost.
  \item The recall share $A/R$, for prediction error and for the safety margin.
  \item Rates of falsely created and missed contexts.
  \item Predicted versus observed values for Propositions~\ref{prop:separation} and~\ref{prop:latency}.
\end{itemize}

%=====================================================================
\section{Risks}
%=====================================================================

\begin{itemize}
  \item \textbf{Contact noise creates false contexts.} This is the most likely problem. The
  state-dependent noise model and a minimum number of steps before a new memory is created address it.
  We report the rate of false contexts as a main metric.
  \item \textbf{Cues are weak.} We vary cue reliability. Proposition~\ref{prop:separation} predicts how
  much a weak cue should still help, so weak cues yield a result rather than a failure.
  \item \textbf{Factors interact.} Mass and friction are expected to interact. Keeping memories for
  seen combinations limits the damage, and the test on unseen combinations measures how large the
  problem is.
  \item \textbf{A change falls outside the correction directions.} The safety margin is measured in the
  full latent space, so safety does not depend on the chosen directions. For prediction, we monitor the
  error outside these directions and add directions for a new memory when this error grows.
  \item \textbf{Latent safety differs from physical safety.} Proposition~\ref{prop:safety} concerns the
  latent condition. We report physical violations separately. Whether $h$ matches the real constraint
  is outside the scope of this proposal~\cite{lutkus2025latent}.
  \item \textbf{Competing work.} AdaJEPA and Sandwich-Residuals name persistence as future work, and
  the latent ACI filters~\cite{huriot2026safe,nath2026pixels} appeared this year. Our distinguishing
  points are acting before contact, combining factors, the first-contact analysis and margins stored
  per memory. A workshop paper by February 2027 establishes priority.
  \item \textbf{The motor-learning patterns add nothing.} Each pattern is tied to a numerical
  prediction and is removed if it does not separate the methods.
\end{itemize}

%=====================================================================
\section{Timeline}
%=====================================================================

\begin{center}
\small
\begin{tabular}{lp{0.45\linewidth}p{0.3\linewidth}}
\toprule
When & Work & Done when \\
\midrule
Oct 2026 & Simulated pushing and puck strike; core of \method{}; gradient baselines & H1, H2 and H4 visible; stopping criteria checked \\
Nov 2026 & Factor memories; safety filter with per-memory margins; recall/learning measurement & H5 and H6 in simulation; Propositions 1--3 compared with simulation \\
Dec 2026--Jan 2027 & AdaJEPA tasks with all baselines; real-robot setup (mats, containers, puck, edge) & Condition sequences running in PushT and PointMaze \\
Feb 2027 & Workshop paper: simulation, AdaJEPA tasks, theory & Submitted \\
Mar--Apr 2027 & Real-robot single-contact tasks; long deployment ($\ge$1000 episodes); safe probing & H1--H8 tested on hardware \\
Spring 2027 & Full paper for CoRL 2027 (deadline to be confirmed when announced) & Submitted \\
\bottomrule
\end{tabular}
\end{center}

\begingroup
\small
\setlength{\bibsep}{1pt}
\bibliographystyle{plainnat}
\bibliography{references}
\endgroup

\end{document}
